% ======================================================================
%  KAVRAM: Prompt & System Prompt  (Tier 1)
% ======================================================================

\concept{Interface · Arayüz}{Prompt \&\\System Prompt}{Modeli değiştiremezsin; değiştirebildiğin tek yüzey gönderdiğin metindir}

% ----------------------------------------------------------------------
\stage{1}{Problem}
% ----------------------------------------------------------------------
Modelin ağırlıklarına dokunamazsın. Onu yeniden yazamaz, içine \texttt{if}
koyamazsın. Elindeki tek kontrol yüzeyi, her çağrıda gönderdiğin \emph{metindir}.
Bu metin muğlaksa, çıktı da muğlak ve tutarsız olur.

\keyline{Prompt bir sohbet değil, bir arayüzdür. Modeli ``programlamanın'' tek yolu odur — ve çoğu ekip onu gelişigüzel yazar.}

İki ekip aynı modeli kullanır; biri tutarlı, ``akıllı'' cevaplar alır, diğeri
rastgele sonuçlarla boğuşur. Fark modelde değil, \textbf{prompt mühendisliğindedir}:
rolü, kuralları ve formatı nereye, nasıl koydukların.

% ----------------------------------------------------------------------
\stage{2}{Neden Var}
% ----------------------------------------------------------------------
LLM durumsuzdur (\en{stateless}): çağrılar arası hiçbir şey hatırlamaz. Her
çağrıda dünyayı sıfırdan, yalnızca o an gönderdiğin metinden kurar. Dolayısıyla:

\begin{itemize}
  \item \textbf{Kalıcı} olması gereken her şey (rol, kurallar, ton, format) her çağrıda yeniden söylenmelidir → \term{system prompt}.
  \item \textbf{O ana özel} istek ayrı tutulmalıdır → \term{user prompt}.
\end{itemize}

Bu ayrım, modele ``kim olduğunu'' ve ``şu an ne istendiğini'' karıştırmadan
söylemeni sağlar.

% ----------------------------------------------------------------------
\stage{3}{Tarihçe}
% ----------------------------------------------------------------------
\begin{timeline}
\yr{2020} completion: ``metni devam ettir'', rol yok \tlsep \yr{2022} instruction tuning: prompt bir komuta dönüşür \tlsep \yr{2023} chat format: \texttt{system/user/assistant} ayrışır \tlsep \yr{2024–26} yapılandırılmış prompt: rol · örnek · şema · araç tanımı tek arayüzde
\end{timeline}

% ----------------------------------------------------------------------
\stage{4}{Tanım}
% ----------------------------------------------------------------------
\begin{problembox}
\textbf{Prompt:} Modele bir çağrıda gönderilen \emph{tüm} girdi.

\smallskip
\textbf{System Prompt:} Rolü, kuralları, tonu ve çıktı formatını tanımlayan,
yüksek öncelikli ve çağrılar boyunca \emph{sabit} kalan talimat katmanı.

\smallskip
\textbf{User Prompt:} O turun özel isteği/sorusu.
\end{problembox}

Modelin gördüğü nihai prompt çoğu zaman bu ikisinin \emph{artı} araya eklenen
bağlamın (geçmiş + getirilen bilgi) birleşimidir.

% ----------------------------------------------------------------------
\stage{5}{Zihinsel Model}
% ----------------------------------------------------------------------
\begin{mentalbox}
System prompt, yeni bir çalışana ilk gün verdiğin \textbf{görev tanımı + ev
kuralları}dır: ``Sen bir hukuk asistanısın, asla tavsiye verme, cevapları madde
madde yaz.'' User prompt ise \textbf{bugünkü somut iş}: ``Şu sözleşmeyi özetle.''

\smallskip
Kritik fark: bu çalışanın \emph{hafızası yok}. Her sabah işe geldiğinde görev
tanımını sıfırdan okur. O yüzden kuralları her çağrıda yeniden vermen gerekir ---
söylemezsen, bilmez.
\end{mentalbox}

% ----------------------------------------------------------------------
\stage{6}{Mimari Diyagram}
% ----------------------------------------------------------------------
Bir prompt'un anatomisi: katmanlar birleşip modele tek girdi olur. Her çağrıda
baştan okunur.

\vspace{4pt}
\begin{center}
\begin{tikzpicture}[
   font=\sffamily\scriptsize,
   layer/.style={rounded corners=2pt, minimum width=58mm, minimum height=11mm,
                 align=left, inner sep=5pt, text width=52mm},
   flow/.style={-{Stealth[length=2mm]}, accent, line width=0.9pt}]

  \node[layer, fill=accentdk, text=white, draw=accentdk, line width=0.9pt]
       (sys) {\textbf{System Prompt}\\[1pt]\footnotesize rol · kurallar · ton · format \,(sabit)};
  \node[layer, fill=wash, draw=accent, line width=0.8pt, below=2mm of sys]
       (ctx) {\textbf{Context}\\[1pt]\footnotesize geçmiş + getirilen bilgi \,(dinamik)};
  \node[layer, fill=white, draw=softgray, line width=0.8pt, below=2mm of ctx]
       (usr) {\textbf{User Prompt}\\[1pt]\footnotesize bu turun isteği};

  \node[draw=hairline, line width=1pt, rounded corners=3pt,
        fit=(sys)(ctx)(usr), inner sep=4mm] (box) {};
  \node[anchor=south west, font=\sffamily\tiny\bfseries, text=softgray]
        at (box.north west) {NİHAİ PROMPT \,(tek girdi)};

  \node[draw=accentdk, fill=accent!12, rounded corners=2pt, right=12mm of ctx,
        minimum width=18mm, minimum height=12mm, align=center] (model) {LLM};
  \node[draw=good, fill=good!10, rounded corners=2pt, right=8mm of model,
        minimum width=18mm, minimum height=10mm, align=center] (out) {Çıktı};

  \draw[flow] (box.east) -- (model);
  \draw[flow] (model) -- (out);
\end{tikzpicture}
\end{center}

\keyline{System en üstte ve en öncelikli; user en altta ve en özel. Araya giren her şey ``context engineering''in konusudur.}

% ----------------------------------------------------------------------
\stage{7}{Karar Çerçevesi}
% ----------------------------------------------------------------------
``Bu bilgi nereye gitmeli --- system'e mi, user'a mı, context'e mi?'' Tek soru
karar verir: \emph{ne kadar sık değişiyor?}

\vspace{2pt}
\begin{center}
\begin{tikzpicture}[
   font=\sffamily\scriptsize, node distance=6mm,
   q/.style={draw=accent, line width=0.8pt, fill=wash, rounded corners=2pt,
             align=center, inner sep=5pt, minimum width=40mm},
   leaf/.style={draw=accentdk, line width=0.8pt, rounded corners=2pt,
                fill=accent!12, align=center, inner sep=4pt, minimum width=30mm},
   ed/.style={-{Stealth[length=2mm]}, softgray, line width=0.7pt},
   lab/.style={font=\sffamily\bfseries\scriptsize, text=accentdk}]

  \node[q] (q) {Bu bilgi ne sıklıkla değişir?};
  \node[leaf, below=14mm of q] (mid) {Her istekte\\\textbf{→ User Prompt}};
  \node[leaf, left=10mm of mid] (sys) {Hiç (sabit kural/rol)\\\textbf{→ System Prompt}};
  \node[leaf, right=10mm of mid] (ctx) {Duruma göre (getirilen)\\\textbf{→ Context}};

  \draw[ed] (q.south) -- (mid.north);
  \draw[ed] (q.south) -- (sys.north);
  \draw[ed] (q.south) -- (ctx.north);
\end{tikzpicture}
\end{center}

\keyline{Kural: Sabit olanı system'e koy ki her çağrıda tutarlı olsun; değişeni user/context'te tut ki esnek kalsın.}

% ----------------------------------------------------------------------
\stage{8}{Üretim Senaryosu}
% ----------------------------------------------------------------------
\begin{prodbox}[Üretim Senaryosu · Tutarsız Destek Botu]
\textbf{Bağlam.} Bir destek botu bazen madde madde, bazen paragraf; bazen Türkçe,
bazen İngilizce; bazen tavsiye veriyor (yasak olmasına rağmen). Kurallar her
\emph{user} mesajının içine elle yazılıyor ve sık unutuluyor.

\smallskip
\textbf{Çözüm.} Tüm kalıcı kurallar (rol, dil, ``asla tıbbi tavsiye verme'',
``en fazla 5 madde'', çıktı formatı) tek bir \textbf{system prompt}'a taşınır.
User prompt yalnızca müşterinin sorusunu içerir. Format için yapılandırılmış
çıktı (JSON) eklenir.

\smallskip
\textbf{Sonuç.} Tutarlılık dramatik artar; kural değişikliği artık 30 farklı yerde
değil, tek bir versiyonlanmış system prompt'ta yapılır.
\end{prodbox}

% ----------------------------------------------------------------------
\stage{9}{Anti-pattern}
% ----------------------------------------------------------------------
\begin{antibox}[Anti-pattern · ``Prompt as Business Logic'']
Kritik iş mantığını --- fiyatlandırma kuralları, yetki kontrolleri, dallanmalar ---
yalnızca uzun bir düzyazı prompt'un içine gömmek. Prompt kırılgandır: tek kelime
değişikliği davranışı sessizce bozar, test edilmesi zordur, denetlenemez.

\smallskip
\textbf{Belirti:} ``Prompt'a bir cümle ekledik, alakasız bir şey bozuldu.''\\
\textbf{Kök neden:} Deterministik olması gereken mantık, olasılıksal düzyazıya emanet edilmiş.\\
\textbf{Düzeltme:} Sert kuralları koda/guardrail'e al; prompt'u yönlendirme için kullan. Prompt'u kod gibi \textbf{versiyonla ve test et}.
\end{antibox}

% ----------------------------------------------------------------------
\stage{10}{Alternatifler \& Ödünleşimler}
% ----------------------------------------------------------------------
\vspace{2pt}
\begin{tabularx}{\linewidth}{@{}>{\raggedright\arraybackslash}p{30mm} L L@{}}
\arrayrulecolor{hairline}\toprule
{\sffamily\bfseries\small Yaklaşım} &
{\sffamily\bfseries\small Güçlü olduğu yer} &
{\sffamily\bfseries\small Ödünleşim} \\
\midrule
\textbf{Düz talimat (prompt)} & Hızlı, esnek, kod gerektirmez &
Tutarlılık prompt kalitesine bağlı \\
\addlinespace[3pt]
\textbf{Few-shot örnek} & Format/üslubu örnekle göstermek &
Token yer; örnekler bağlamı şişirir \\
\addlinespace[3pt]
\textbf{Structured output} & Format garantisi gerektiğinde &
Esnekliği azaltır; şema bakımı \\
\addlinespace[3pt]
\textbf{Fine-tuning} & Sabit üslubu kalıcı gömmek &
Pahalı; her değişiklik yeniden eğitim \\
\bottomrule
\end{tabularx}

% ----------------------------------------------------------------------
\stage{11}{İlişkili Kavramlar}
% ----------------------------------------------------------------------
\begin{itemize}
  \item \term{Context Engineering} --- prompt'a neyin, hangi sırayla gireceğini yönetir.
  \item \term{Few-shot} --- prompt içine konan örneklerle yönlendirme.
  \item \term{Structured Output} --- çıktı formatını şemayla garanti altına alma.
  \item \term{Temperature} --- aynı prompt'un ne kadar çeşitli yorumlanacağı.
  \item \term{Prompt Versioning} --- prompt'u kod gibi izleyip geri alabilme.
\end{itemize}

% ----------------------------------------------------------------------
\stage{12}{Özet}
% ----------------------------------------------------------------------
\begin{itemize}
  \item Prompt, modeli yönlendirebildiğin \textbf{tek kontrol yüzeyidir}.
  \item Model durumsuzdur: kalıcı kurallar her çağrıda system prompt'ta yeniden verilmelidir.
  \item Sabit olanı system'e, isteğe özel olanı user'a, dinamik bilgiyi context'e koy.
  \item Sert iş mantığını prompt'a gömme --- onu koda/guardrail'e al.
  \item Prompt'u kod gibi gör: versiyonla, test et, gözden geçir.
\end{itemize}

% ----------------------------------------------------------------------
\stage{13}{Mühendislik İlkesi}
% ----------------------------------------------------------------------
\begin{principlebox}
\textbf{İlke \#2 — Prompt koddur: istikrarlı olanı system'e koy, versiyonla, test et.}\\[4pt]
\normalsize Modeli değiştiremezsin; yalnızca ona ne söylediğini değiştirebilirsin.
O hâlde bu yüzeye bir sohbet gibi değil, bir arayüz gibi davran.
\end{principlebox}
